% ======================================================================
% Beamer presentation - generated by KherveSlide.
% Edit freely; the source is regenerated when you change a slide.
% ======================================================================
\documentclass[aspectratio=169]{beamer}

% --- theme & packages ---
\usetheme{default}
\definecolor{ksth0}{HTML}{003E74}
\setbeamercolor{structure}{fg=ksth0}
\definecolor{ksth1}{HTML}{212121}
\setbeamercolor{normal text}{fg=ksth1}
\definecolor{ksth2}{HTML}{FFFFFF}
\definecolor{ksth3}{HTML}{003E74}
\setbeamercolor{frametitle}{fg=ksth2,bg=ksth3}
\definecolor{ksth4}{HTML}{FFFFFF}
\setbeamercolor{title}{fg=ksth4}
\definecolor{ksth5}{HTML}{D1DCE6}
\setbeamercolor{block title}{bg=ksth5}
\definecolor{ksth6}{HTML}{FFFFFF}
\definecolor{ksth7}{HTML}{003E74}
\setbeamercolor{ks footer bar}{fg=ksth6,bg=ksth7}
\setbeamertemplate{footline}{%
\leavevmode\hbox{\begin{beamercolorbox}[wd=\paperwidth,ht=2.6ex,dp=1.1ex,leftskip=1.5ex,rightskip=1.5ex]{ks footer bar}%
\usebeamerfont{author in head/foot}\insertshortauthor\hfill\insertshorttitle\hfill\insertframenumber\,/\,\inserttotalframenumber
\end{beamercolorbox}}\vskip0pt}
\usepackage{lmodern}
\usepackage[version=4]{mhchem}
\usepackage{tikz}
\usetikzlibrary{arrows.meta}
\usepackage{colortbl}
\definecolor{ksTblHead}{HTML}{FCE4D6}
\definecolor{ksTblHeadFg}{HTML}{C55A11}
\definecolor{ksTblRule}{HTML}{F4B183}
\definecolor{ksTblCap}{HTML}{808080}
\usepackage[absolute,overlay]{textpos}
\usepackage{graphicx}
\setlength{\TPHorizModule}{1\paperwidth}
\setlength{\TPVertModule}{1\paperheight}
\textblockorigin{0\paperwidth}{0\paperheight}
\setbeamerfont{itemize/enumerate subbody}{size=\relax}
\setbeamerfont{itemize/enumerate subsubbody}{size=\relax}
\title{Oxygen vacancies in TiO2 thin films}
\author{A. Researcher}

\begin{document}

% ======================================================================
% Slide 1
% ======================================================================
\begin{frame}[t]

  % text box (free-positioned)
\begin{textblock}{0.84}(0.08,0.16)
{\centering\fontsize{32}{38}\selectfont \textcolor[HTML]{1A1A1A}{\textbf{Oxygen vacancies in TiO$_2$ thin films}}\par}
\end{textblock}

  % line / arrow (free-positioned)
\definecolor{ksline0}{HTML}{003E74}
\begin{textblock}{0.24}(0.38,0.52)
\begin{tikzpicture}
\useasboundingbox (0,0) rectangle (\linewidth,-0\TPVertModule);
\draw[line width=2.5pt,color=ksline0] (0\linewidth,-0\TPVertModule) -- (1\linewidth,-0\TPVertModule);
\end{tikzpicture}
\end{textblock}

  % text box (free-positioned)
\begin{textblock}{0.84}(0.08,0.57)
{\centering\fontsize{17}{20}\selectfont \textcolor[HTML]{555555}{Sputter deposition, XRD, XPS and optical band gap}\par}
\end{textblock}

  % text box (free-positioned)
\begin{textblock}{0.84}(0.08,0.76)
{\centering\fontsize{13}{16}\selectfont \textcolor[HTML]{555555}{A. Researcher  ·  Surface Analysis Group  ·  \today}\par}
\end{textblock}
\end{frame}

% ======================================================================
% Slide 2 - Samples and workflow
% ======================================================================
\begin{frame}[t]
  \frametitle{Samples and workflow}

  % shape (free-positioned)
\definecolor{ksfill1}{HTML}{003E74}%
\begin{textblock}{0.2}(0.04,0.24)
\begin{tikzpicture}
\useasboundingbox (0,0) rectangle (\linewidth,-0.16\TPVertModule);
\path[fill=ksfill1] (0\linewidth,-0\TPVertModule) -- (0.72\linewidth,-0\TPVertModule) -- (1\linewidth,-0.08\TPVertModule) -- (0.72\linewidth,-0.16\TPVertModule) -- (0\linewidth,-0.16\TPVertModule) -- (0.28\linewidth,-0.08\TPVertModule) -- cycle;
\end{tikzpicture}
\end{textblock}

  % text box (free-positioned)
\begin{textblock}{0.16}(0.06,0.301)
{\centering\fontsize{11}{13}\selectfont \textcolor[HTML]{FFFFFF}{\textbf{Sputter}}\par}
\end{textblock}

  % shape (free-positioned)
\definecolor{ksfill2}{HTML}{0091D4}%
\begin{textblock}{0.2}(0.23,0.24)
\begin{tikzpicture}
\useasboundingbox (0,0) rectangle (\linewidth,-0.16\TPVertModule);
\path[fill=ksfill2] (0\linewidth,-0\TPVertModule) -- (0.72\linewidth,-0\TPVertModule) -- (1\linewidth,-0.08\TPVertModule) -- (0.72\linewidth,-0.16\TPVertModule) -- (0\linewidth,-0.16\TPVertModule) -- (0.28\linewidth,-0.08\TPVertModule) -- cycle;
\end{tikzpicture}
\end{textblock}

  % text box (free-positioned)
\begin{textblock}{0.16}(0.25,0.301)
{\centering\fontsize{11}{13}\selectfont \textcolor[HTML]{FFFFFF}{\textbf{Anneal}}\par}
\end{textblock}

  % shape (free-positioned)
\definecolor{ksfill3}{HTML}{2E8B57}%
\begin{textblock}{0.2}(0.42,0.24)
\begin{tikzpicture}
\useasboundingbox (0,0) rectangle (\linewidth,-0.16\TPVertModule);
\path[fill=ksfill3] (0\linewidth,-0\TPVertModule) -- (0.72\linewidth,-0\TPVertModule) -- (1\linewidth,-0.08\TPVertModule) -- (0.72\linewidth,-0.16\TPVertModule) -- (0\linewidth,-0.16\TPVertModule) -- (0.28\linewidth,-0.08\TPVertModule) -- cycle;
\end{tikzpicture}
\end{textblock}

  % text box (free-positioned)
\begin{textblock}{0.16}(0.44,0.301)
{\centering\fontsize{11}{13}\selectfont \textcolor[HTML]{FFFFFF}{\textbf{XRD}}\par}
\end{textblock}

  % shape (free-positioned)
\definecolor{ksfill4}{HTML}{C8102E}%
\begin{textblock}{0.2}(0.61,0.24)
\begin{tikzpicture}
\useasboundingbox (0,0) rectangle (\linewidth,-0.16\TPVertModule);
\path[fill=ksfill4] (0\linewidth,-0\TPVertModule) -- (0.72\linewidth,-0\TPVertModule) -- (1\linewidth,-0.08\TPVertModule) -- (0.72\linewidth,-0.16\TPVertModule) -- (0\linewidth,-0.16\TPVertModule) -- (0.28\linewidth,-0.08\TPVertModule) -- cycle;
\end{tikzpicture}
\end{textblock}

  % text box (free-positioned)
\begin{textblock}{0.16}(0.63,0.301)
{\centering\fontsize{11}{13}\selectfont \textcolor[HTML]{FFFFFF}{\textbf{XPS}}\par}
\end{textblock}

  % shape (free-positioned)
\definecolor{ksfill5}{HTML}{B8860B}%
\begin{textblock}{0.2}(0.8,0.24)
\begin{tikzpicture}
\useasboundingbox (0,0) rectangle (\linewidth,-0.16\TPVertModule);
\path[fill=ksfill5] (0\linewidth,-0\TPVertModule) -- (0.72\linewidth,-0\TPVertModule) -- (1\linewidth,-0.08\TPVertModule) -- (0.72\linewidth,-0.16\TPVertModule) -- (0\linewidth,-0.16\TPVertModule) -- (0.28\linewidth,-0.08\TPVertModule) -- cycle;
\end{tikzpicture}
\end{textblock}

  % text box (free-positioned)
\begin{textblock}{0.16}(0.82,0.301)
{\centering\fontsize{11}{13}\selectfont \textcolor[HTML]{FFFFFF}{\textbf{UV–Vis}}\par}
\end{textblock}

  % text box (free-positioned)
\begin{textblock}{0.56}(0.06,0.5)
{\raggedright\fontsize{15}{18}\selectfont \begin{itemize}
\item Reactive DC sputtering, Ar/O$_2$ = 4:1, 300\,W
\item Anneal 450\,°C, 2\,h, in air or vacuum
\item Films 200\,nm on Si and quartz
\end{itemize}\par}
\end{textblock}

  % text box (free-positioned)
\begin{textblock}{0.3}(0.65,0.5)
{\begin{block}{Hypothesis}\raggedright\fontsize{13}{16}\selectfont Vacuum anneal removes lattice oxygen: $\ce{TiO2 -> TiO_{2-x} + x/2 O2}$\end{block}\par}
\end{textblock}
\end{frame}

% ======================================================================
% Slide 3 - Structure: anatase only
% ======================================================================
\begin{frame}[t]
  \frametitle{Structure: anatase only}

  % picture (free-positioned)
\begin{textblock}{0.56}(0.03,0.18)
\includegraphics[width=0.56\paperwidth,height=0.66\paperheight,keepaspectratio]{media/example_xrd.png}
\end{textblock}

  % text box (free-positioned)
\begin{textblock}{0.34}(0.62,0.22)
{\raggedright\fontsize{14}{17}\selectfont Crystallite size (Scherrer):\par}
\end{textblock}

  % text box (free-positioned)
\begin{textblock}{0.34}(0.62,0.28)
{\centering\fontsize{18}{22}\selectfont \[D = \frac{K\lambda}{\beta\cos\theta}\]\par}
\end{textblock}

  % table (free-positioned)
\begin{textblock}{0.26}(0.66,0.52)
{\definecolor{ksTH003E74}{HTML}{003E74}\definecolor{ksTFFFFFFF}{HTML}{FFFFFF}\definecolor{ksTRC8C8C8}{HTML}{C8C8C8}\definecolor{ksTSF5F5F5}{HTML}{F5F5F5}\setlength{\arrayrulewidth}{0.8pt}\arrayrulecolor{ksTRC8C8C8} \fontsize{13}{16}\selectfont \begin{tabular}{|c|c|}
\hline
  \rowcolor{ksTH003E74}\textcolor{ksTFFFFFFF}{\textbf{Anneal}} & \textcolor{ksTFFFFFFF}{\textbf{$D$ (nm)}} \\
\hline
  air & 18 \\
\hline
  vacuum & 21 \\
\hline
\end{tabular}}
\end{textblock}

  % text box (free-positioned)
\begin{textblock}{0.34}(0.62,0.77)
{\raggedright\fontsize{12}{14}\selectfont \textcolor[HTML]{555555}{\textit{No rutile or brookite detected.}}\par}
\end{textblock}
\end{frame}

% ======================================================================
% Slide 4 - Chemistry: Ti 2p XPS
% ======================================================================
\begin{frame}[t]
  \frametitle{Chemistry: Ti 2p XPS}

  % picture (free-positioned)
\begin{textblock}{0.56}(0.03,0.18)
\includegraphics[width=0.56\paperwidth,height=0.66\paperheight,keepaspectratio]{media/example_xps.png}
\end{textblock}

  % text box (free-positioned)
\begin{textblock}{0.36}(0.61,0.2)
{\raggedright\fontsize{13}{16}\selectfont \begin{itemize}
\item Doublet split 5.7\,eV, area ratio 2:1
\item Ti$^{3+}$ shoulder at 457.1\,eV
\item Shirley background, pseudo-Voigt
\end{itemize}\par}
\end{textblock}

  % table (free-positioned)
\begin{textblock}{0.34}(0.62,0.64)
{\definecolor{ksTH003E74}{HTML}{003E74}\definecolor{ksTFFFFFFF}{HTML}{FFFFFF}\definecolor{ksTRC8C8C8}{HTML}{C8C8C8}\definecolor{ksTSF5F5F5}{HTML}{F5F5F5}\setlength{\arrayrulewidth}{0.8pt}\arrayrulecolor{ksTRC8C8C8} \fontsize{12}{14}\selectfont \begin{tabular}{|c|c|c|}
\hline
  \rowcolor{ksTH003E74}\textcolor{ksTFFFFFFF}{\textbf{}} & \textcolor{ksTFFFFFFF}{\textbf{Ti$^{4+}$}} & \textcolor{ksTFFFFFFF}{\textbf{Ti$^{3+}$}} \\
\hline
  air & 100\,\% & — \\
\hline
  vacuum & 89\,\% & 11\,\% \\
\hline
\end{tabular}}
\end{textblock}
\end{frame}

% ======================================================================
% Slide 5 - Optical band gap
% ======================================================================
\begin{frame}[t]
  \frametitle{Optical band gap}

  % picture (free-positioned)
\begin{textblock}{0.48}(0.04,0.18)
\includegraphics[width=0.48\paperwidth,height=0.66\paperheight,keepaspectratio]{media/example_tauc.png}
\end{textblock}

  % text box (free-positioned)
\begin{textblock}{0.4}(0.56,0.22)
{\raggedright\fontsize{14}{17}\selectfont Tauc relation (indirect gap, $n=2$):\par}
\end{textblock}

  % text box (free-positioned)
\begin{textblock}{0.42}(0.55,0.3)
{\centering\fontsize{20}{24}\selectfont \[(\alpha h\nu)^{1/n} = A\,(h\nu - E_g)\]\par}
\end{textblock}

  % text box (free-positioned)
\begin{textblock}{0.4}(0.56,0.52)
{\begin{alertblock}{Observation}\raggedright\fontsize{13}{16}\selectfont Vacancies add states 0.8\,eV below the conduction band → visible absorption, grey-blue films.\end{alertblock}\par}
\end{textblock}
\end{frame}

% ======================================================================
% Slide 6 - Conclusions
% ======================================================================
\begin{frame}[t]
  \frametitle{Conclusions}

  % shape (free-positioned)
\definecolor{ksfill6}{HTML}{003E74}%
\begin{textblock}{0.24}(0.08,0.22)
\begin{tikzpicture}
\useasboundingbox (0,0) rectangle (\linewidth,-0.22\TPVertModule);
\path[fill=ksfill6, rounded corners=6pt] (0\linewidth,-0\TPVertModule) rectangle (1\linewidth,-0.22\TPVertModule);
\end{tikzpicture}
\end{textblock}

  % text box (free-positioned)
\begin{textblock}{0.24}(0.08,0.261)
{\centering\fontsize{26}{31}\selectfont \textcolor[HTML]{FFFFFF}{\textbf{11\,\%}}\par}
\end{textblock}

  % text box (free-positioned)
\begin{textblock}{0.24}(0.08,0.358)
{\centering\fontsize{12}{14}\selectfont \textcolor[HTML]{FFFFFF}{Ti$^{3+}$ after vacuum}\par}
\end{textblock}

  % shape (free-positioned)
\definecolor{ksfill7}{HTML}{0091D4}%
\begin{textblock}{0.24}(0.38,0.22)
\begin{tikzpicture}
\useasboundingbox (0,0) rectangle (\linewidth,-0.22\TPVertModule);
\path[fill=ksfill7, rounded corners=6pt] (0\linewidth,-0\TPVertModule) rectangle (1\linewidth,-0.22\TPVertModule);
\end{tikzpicture}
\end{textblock}

  % text box (free-positioned)
\begin{textblock}{0.24}(0.38,0.261)
{\centering\fontsize{26}{31}\selectfont \textcolor[HTML]{FFFFFF}{\textbf{3.22 eV}}\par}
\end{textblock}

  % text box (free-positioned)
\begin{textblock}{0.24}(0.38,0.358)
{\centering\fontsize{12}{14}\selectfont \textcolor[HTML]{FFFFFF}{band gap}\par}
\end{textblock}

  % shape (free-positioned)
\definecolor{ksfill8}{HTML}{2E8B57}%
\begin{textblock}{0.24}(0.68,0.22)
\begin{tikzpicture}
\useasboundingbox (0,0) rectangle (\linewidth,-0.22\TPVertModule);
\path[fill=ksfill8, rounded corners=6pt] (0\linewidth,-0\TPVertModule) rectangle (1\linewidth,-0.22\TPVertModule);
\end{tikzpicture}
\end{textblock}

  % text box (free-positioned)
\begin{textblock}{0.24}(0.68,0.261)
{\centering\fontsize{26}{31}\selectfont \textcolor[HTML]{FFFFFF}{\textbf{21 nm}}\par}
\end{textblock}

  % text box (free-positioned)
\begin{textblock}{0.24}(0.68,0.358)
{\centering\fontsize{12}{14}\selectfont \textcolor[HTML]{FFFFFF}{crystallites}\par}
\end{textblock}

  % text box (free-positioned)
\begin{textblock}{0.84}(0.08,0.52)
{\raggedright\fontsize{16}{19}\selectfont \begin{itemize}
\item Vacuum annealing creates Ti$^{3+}$ without a phase change
\item XPS quantifies it; XRD and UV–Vis confirm
\item Next: photocatalysis against vacancy density
\end{itemize}\par}
\end{textblock}

  % text box (free-positioned)
\begin{textblock}{0.8}(0.1,0.86)
{\centering\fontsize{10}{12}\selectfont \textcolor[HTML]{777777}{Spectra fitted in KherveFitting  ·  slides in KherveSlide}\par}
\end{textblock}
\end{frame}

\end{document}
